%%==========================================================================
%%  saus-exam-demo.tex
%%  A final-term paper showing every element of saus-exam.cls: the heading,
%%  the examiner's grid, instructions, sections, multiple-choice items,
%%  questions in parts, a table, an equation, code, answer space, bonus
%%  marks, model answers and the marks/CLO distribution table.
%%  Compile with:  latexmk saus-exam-demo.tex   (LuaLaTeX via .latexmkrc)
%%
%%  For the answer key from this same file, put these two lines in a file of
%%  their own and compile that instead:
%%    \PassOptionsToClass{solutions=1}{saus-exam}
%%    %%==========================================================================
%%  saus-exam-demo.tex
%%  A final-term paper showing every element of saus-exam.cls: the heading,
%%  the examiner's grid, instructions, sections, multiple-choice items,
%%  questions in parts, a table, an equation, code, answer space, bonus
%%  marks, model answers and the marks/CLO distribution table.
%%  Compile with:  latexmk saus-exam-demo.tex   (LuaLaTeX via .latexmkrc)
%%
%%  For the answer key from this same file, put these two lines in a file of
%%  their own and compile that instead:
%%    \PassOptionsToClass{solutions=1}{saus-exam}
%%    %%==========================================================================
%%  saus-exam-demo.tex
%%  A final-term paper showing every element of saus-exam.cls: the heading,
%%  the examiner's grid, instructions, sections, multiple-choice items,
%%  questions in parts, a table, an equation, code, answer space, bonus
%%  marks, model answers and the marks/CLO distribution table.
%%  Compile with:  latexmk saus-exam-demo.tex   (LuaLaTeX via .latexmkrc)
%%
%%  For the answer key from this same file, put these two lines in a file of
%%  their own and compile that instead:
%%    \PassOptionsToClass{solutions=1}{saus-exam}
%%    %%==========================================================================
%%  saus-exam-demo.tex
%%  A final-term paper showing every element of saus-exam.cls: the heading,
%%  the examiner's grid, instructions, sections, multiple-choice items,
%%  questions in parts, a table, an equation, code, answer space, bonus
%%  marks, model answers and the marks/CLO distribution table.
%%  Compile with:  latexmk saus-exam-demo.tex   (LuaLaTeX via .latexmkrc)
%%
%%  For the answer key from this same file, put these two lines in a file of
%%  their own and compile that instead:
%%    \PassOptionsToClass{solutions=1}{saus-exam}
%%    \input{saus-exam-demo}
%%==========================================================================
% type = mid | final | sessional | quiz | makeup | supplementary | practical
% solutions = 1 prints the model answers and marks the correct choices
\documentclass[type=final]{saus-exam}

%% ---- The paper -------------------------------------------------------------
\sausdepartment{Department of Computer Science}
\examcourse{CSC-201}{Data Structures and Algorithms}
\examprogramme{BS Computer Science}
\examsemester{III}
\examsession{Spring 2026}
\examdate{16 June 2026}
\examduration{3 hours}
\exammarks{50}
\examcredits{3 (3-0)}
\examinfo{Venue}{Examination Hall 2}

\begin{document}
\makeexamhead

\begin{instructions}
  \item Attempt \textbf{all} questions. Section A is answered on this paper;
        Sections B and C in the answer book provided.
  \item The marks for each question are shown in the margin, with the course
        learning outcome it assesses.
  \item Programmable calculators, mobile telephones and smart watches are not
        allowed in the examination hall.
  \item Write your roll number on every sheet of the answer book.
\end{instructions}

\sausmarksgrid

\examsection[10 marks]{Section A: Multiple choice}

\question[marks=2, clo=1, bloom=Remember, topic=Complexity]
What is the worst-case time complexity of binary search on a sorted array
of $n$ elements?
\begin{choices}[cols=4]
  \choice $O(1)$
  \correctchoice $O(\log n)$
  \choice $O(n)$
  \choice $O(n\log n)$
\end{choices}
\begin{solution}
  Every comparison halves the interval that may still contain the key, so at
  most $\lfloor\log_2 n\rfloor+1$ comparisons are made.
\end{solution}

\question[marks=2, clo=1, bloom=Understand, topic=Stacks and queues]
Which structure serves its items in last-in, first-out order?
\begin{choices}
  \choice Queue
  \correctchoice Stack
  \choice Priority queue
  \choice Circular buffer
\end{choices}

\question[marks=2, clo=1, bloom=Understand, topic=Trees]
A binary search tree holding $n$ keys has height $h$. Searching it takes
time proportional to:
\begin{choices}
  \choice $n$ in every case
  \correctchoice $h$, which is $\log n$ only when the tree is balanced
  \choice $\log n$ in every case
  \choice $n\log n$ in the worst case
\end{choices}

\question[marks=2, clo=2, bloom=Apply, topic=Hashing]
Separate chaining resolves a collision by:
\begin{choices}
  \choice probing the next free slot
  \correctchoice keeping a list of the entries that share a slot
  \choice enlarging the table on every insertion
  \choice rehashing with a second function
\end{choices}

\question[marks=2, clo=2, bloom=Apply, topic=Graphs]
Breadth-first search on an unweighted graph finds:
\begin{choices}
  \choice the minimum spanning tree
  \correctchoice the shortest path in edges from the source
  \choice the strongly connected components
  \choice a topological order
\end{choices}

\examsection[20 marks]{Section B: Short questions}

\question[marks=5, clo=2, bloom=Apply, topic=Recursion]
Write a recursive function that returns the height of a binary tree, and
state its running time in terms of the number of nodes.
\begin{solution}
  The height of an empty tree is $-1$; otherwise it is one more than the
  greater of the heights of the two subtrees. Every node is visited once, so
  the function runs in $O(n)$ time and uses $O(h)$ stack space.
\end{solution}

\question[marks=5, clo=2, bloom=Analyse, topic=Sorting]
Complete the table below, and say which of the three you would choose for
an array that is already nearly sorted.
\begin{center}
  \small
  \begin{tabular}{|l|l|l|l|}
    \hline
    \rowcolor{sausmist}
    \textbf{Algorithm}&\textbf{Best}&\textbf{Average}&\textbf{Worst}\\
    \hline
    Insertion sort&$O(n)$&&\\
    \hline
    Merge sort&&$O(n\log n)$&\\
    \hline
    Quicksort&&&\\
    \hline
  \end{tabular}
\end{center}

\question[marks=5, clo=3, bloom=Analyse, topic=Complexity]
The running time of a divide-and-conquer algorithm satisfies
\begin{equation*}
  T(n) = 2\,T\!\left(\frac{n}{2}\right) + \Theta(n), \qquad T(1)=\Theta(1).
\end{equation*}
Solve the recurrence and name an algorithm whose running time it describes.

\question[marks=5, clo=3, bloom=Evaluate, topic=Data structure choice,
          lines=6]
A campus application keeps the 815 enrolled students in memory and looks a
student up by roll number thousands of times a minute, while enrolments are
added only a few times a day. Recommend a data structure and justify the
choice in terms of the operations' costs.

\examsection[20 marks]{Section C: Long questions}

\question[marks=10, clo=3, bloom=Analyse, topic=Sorting]
Consider the array $[38, 27, 43, 3, 9, 82, 10]$.
\qpart[marks=6]
Trace merge sort on it, showing the contents of the array after every merge.
\qpart[marks=4]
Count the comparisons merge sort makes on this input and compare the number
with insertion sort on the same input.

\question[marks=10, clo=4, bloom=Create, topic=Hashing]
A hash table stores the courses a department offers, keyed by course code.
\qpart[marks=4]
Write the insertion routine for separate chaining. One line of code per line
of the answer book:
\begin{center}\small
  \texttt{def insert(table, key, value):}
\end{center}
\qpart[marks=6, space=45mm]
The table has 16 slots and holds 40 courses. Explain what the load factor
does to the cost of a lookup, and describe one way to keep the cost near
constant as the department adds courses. Use the space below for a diagram.

\question[marks=3, bonus, clo=4, bloom=Create, topic=Open question]
State one change to the hash table above that would make its iteration order
predictable, and what it would cost.

\examend

%% ---- For the department's file ---------------------------------------------
%% The distribution table is for the paper setter and the examination office.
%% Delete this page before printing the candidates' copies.
\clearpage
\sausmarkstable

\end{document}

%%==========================================================================
% type = mid | final | sessional | quiz | makeup | supplementary | practical
% solutions = 1 prints the model answers and marks the correct choices
\documentclass[type=final]{saus-exam}

%% ---- The paper -------------------------------------------------------------
\sausdepartment{Department of Computer Science}
\examcourse{CSC-201}{Data Structures and Algorithms}
\examprogramme{BS Computer Science}
\examsemester{III}
\examsession{Spring 2026}
\examdate{16 June 2026}
\examduration{3 hours}
\exammarks{50}
\examcredits{3 (3-0)}
\examinfo{Venue}{Examination Hall 2}

\begin{document}
\makeexamhead

\begin{instructions}
  \item Attempt \textbf{all} questions. Section A is answered on this paper;
        Sections B and C in the answer book provided.
  \item The marks for each question are shown in the margin, with the course
        learning outcome it assesses.
  \item Programmable calculators, mobile telephones and smart watches are not
        allowed in the examination hall.
  \item Write your roll number on every sheet of the answer book.
\end{instructions}

\sausmarksgrid

\examsection[10 marks]{Section A: Multiple choice}

\question[marks=2, clo=1, bloom=Remember, topic=Complexity]
What is the worst-case time complexity of binary search on a sorted array
of $n$ elements?
\begin{choices}[cols=4]
  \choice $O(1)$
  \correctchoice $O(\log n)$
  \choice $O(n)$
  \choice $O(n\log n)$
\end{choices}
\begin{solution}
  Every comparison halves the interval that may still contain the key, so at
  most $\lfloor\log_2 n\rfloor+1$ comparisons are made.
\end{solution}

\question[marks=2, clo=1, bloom=Understand, topic=Stacks and queues]
Which structure serves its items in last-in, first-out order?
\begin{choices}
  \choice Queue
  \correctchoice Stack
  \choice Priority queue
  \choice Circular buffer
\end{choices}

\question[marks=2, clo=1, bloom=Understand, topic=Trees]
A binary search tree holding $n$ keys has height $h$. Searching it takes
time proportional to:
\begin{choices}
  \choice $n$ in every case
  \correctchoice $h$, which is $\log n$ only when the tree is balanced
  \choice $\log n$ in every case
  \choice $n\log n$ in the worst case
\end{choices}

\question[marks=2, clo=2, bloom=Apply, topic=Hashing]
Separate chaining resolves a collision by:
\begin{choices}
  \choice probing the next free slot
  \correctchoice keeping a list of the entries that share a slot
  \choice enlarging the table on every insertion
  \choice rehashing with a second function
\end{choices}

\question[marks=2, clo=2, bloom=Apply, topic=Graphs]
Breadth-first search on an unweighted graph finds:
\begin{choices}
  \choice the minimum spanning tree
  \correctchoice the shortest path in edges from the source
  \choice the strongly connected components
  \choice a topological order
\end{choices}

\examsection[20 marks]{Section B: Short questions}

\question[marks=5, clo=2, bloom=Apply, topic=Recursion]
Write a recursive function that returns the height of a binary tree, and
state its running time in terms of the number of nodes.
\begin{solution}
  The height of an empty tree is $-1$; otherwise it is one more than the
  greater of the heights of the two subtrees. Every node is visited once, so
  the function runs in $O(n)$ time and uses $O(h)$ stack space.
\end{solution}

\question[marks=5, clo=2, bloom=Analyse, topic=Sorting]
Complete the table below, and say which of the three you would choose for
an array that is already nearly sorted.
\begin{center}
  \small
  \begin{tabular}{|l|l|l|l|}
    \hline
    \rowcolor{sausmist}
    \textbf{Algorithm}&\textbf{Best}&\textbf{Average}&\textbf{Worst}\\
    \hline
    Insertion sort&$O(n)$&&\\
    \hline
    Merge sort&&$O(n\log n)$&\\
    \hline
    Quicksort&&&\\
    \hline
  \end{tabular}
\end{center}

\question[marks=5, clo=3, bloom=Analyse, topic=Complexity]
The running time of a divide-and-conquer algorithm satisfies
\begin{equation*}
  T(n) = 2\,T\!\left(\frac{n}{2}\right) + \Theta(n), \qquad T(1)=\Theta(1).
\end{equation*}
Solve the recurrence and name an algorithm whose running time it describes.

\question[marks=5, clo=3, bloom=Evaluate, topic=Data structure choice,
          lines=6]
A campus application keeps the 815 enrolled students in memory and looks a
student up by roll number thousands of times a minute, while enrolments are
added only a few times a day. Recommend a data structure and justify the
choice in terms of the operations' costs.

\examsection[20 marks]{Section C: Long questions}

\question[marks=10, clo=3, bloom=Analyse, topic=Sorting]
Consider the array $[38, 27, 43, 3, 9, 82, 10]$.
\qpart[marks=6]
Trace merge sort on it, showing the contents of the array after every merge.
\qpart[marks=4]
Count the comparisons merge sort makes on this input and compare the number
with insertion sort on the same input.

\question[marks=10, clo=4, bloom=Create, topic=Hashing]
A hash table stores the courses a department offers, keyed by course code.
\qpart[marks=4]
Write the insertion routine for separate chaining. One line of code per line
of the answer book:
\begin{center}\small
  \texttt{def insert(table, key, value):}
\end{center}
\qpart[marks=6, space=45mm]
The table has 16 slots and holds 40 courses. Explain what the load factor
does to the cost of a lookup, and describe one way to keep the cost near
constant as the department adds courses. Use the space below for a diagram.

\question[marks=3, bonus, clo=4, bloom=Create, topic=Open question]
State one change to the hash table above that would make its iteration order
predictable, and what it would cost.

\examend

%% ---- For the department's file ---------------------------------------------
%% The distribution table is for the paper setter and the examination office.
%% Delete this page before printing the candidates' copies.
\clearpage
\sausmarkstable

\end{document}

%%==========================================================================
% type = mid | final | sessional | quiz | makeup | supplementary | practical
% solutions = 1 prints the model answers and marks the correct choices
\documentclass[type=final]{saus-exam}

%% ---- The paper -------------------------------------------------------------
\sausdepartment{Department of Computer Science}
\examcourse{CSC-201}{Data Structures and Algorithms}
\examprogramme{BS Computer Science}
\examsemester{III}
\examsession{Spring 2026}
\examdate{16 June 2026}
\examduration{3 hours}
\exammarks{50}
\examcredits{3 (3-0)}
\examinfo{Venue}{Examination Hall 2}

\begin{document}
\makeexamhead

\begin{instructions}
  \item Attempt \textbf{all} questions. Section A is answered on this paper;
        Sections B and C in the answer book provided.
  \item The marks for each question are shown in the margin, with the course
        learning outcome it assesses.
  \item Programmable calculators, mobile telephones and smart watches are not
        allowed in the examination hall.
  \item Write your roll number on every sheet of the answer book.
\end{instructions}

\sausmarksgrid

\examsection[10 marks]{Section A: Multiple choice}

\question[marks=2, clo=1, bloom=Remember, topic=Complexity]
What is the worst-case time complexity of binary search on a sorted array
of $n$ elements?
\begin{choices}[cols=4]
  \choice $O(1)$
  \correctchoice $O(\log n)$
  \choice $O(n)$
  \choice $O(n\log n)$
\end{choices}
\begin{solution}
  Every comparison halves the interval that may still contain the key, so at
  most $\lfloor\log_2 n\rfloor+1$ comparisons are made.
\end{solution}

\question[marks=2, clo=1, bloom=Understand, topic=Stacks and queues]
Which structure serves its items in last-in, first-out order?
\begin{choices}
  \choice Queue
  \correctchoice Stack
  \choice Priority queue
  \choice Circular buffer
\end{choices}

\question[marks=2, clo=1, bloom=Understand, topic=Trees]
A binary search tree holding $n$ keys has height $h$. Searching it takes
time proportional to:
\begin{choices}
  \choice $n$ in every case
  \correctchoice $h$, which is $\log n$ only when the tree is balanced
  \choice $\log n$ in every case
  \choice $n\log n$ in the worst case
\end{choices}

\question[marks=2, clo=2, bloom=Apply, topic=Hashing]
Separate chaining resolves a collision by:
\begin{choices}
  \choice probing the next free slot
  \correctchoice keeping a list of the entries that share a slot
  \choice enlarging the table on every insertion
  \choice rehashing with a second function
\end{choices}

\question[marks=2, clo=2, bloom=Apply, topic=Graphs]
Breadth-first search on an unweighted graph finds:
\begin{choices}
  \choice the minimum spanning tree
  \correctchoice the shortest path in edges from the source
  \choice the strongly connected components
  \choice a topological order
\end{choices}

\examsection[20 marks]{Section B: Short questions}

\question[marks=5, clo=2, bloom=Apply, topic=Recursion]
Write a recursive function that returns the height of a binary tree, and
state its running time in terms of the number of nodes.
\begin{solution}
  The height of an empty tree is $-1$; otherwise it is one more than the
  greater of the heights of the two subtrees. Every node is visited once, so
  the function runs in $O(n)$ time and uses $O(h)$ stack space.
\end{solution}

\question[marks=5, clo=2, bloom=Analyse, topic=Sorting]
Complete the table below, and say which of the three you would choose for
an array that is already nearly sorted.
\begin{center}
  \small
  \begin{tabular}{|l|l|l|l|}
    \hline
    \rowcolor{sausmist}
    \textbf{Algorithm}&\textbf{Best}&\textbf{Average}&\textbf{Worst}\\
    \hline
    Insertion sort&$O(n)$&&\\
    \hline
    Merge sort&&$O(n\log n)$&\\
    \hline
    Quicksort&&&\\
    \hline
  \end{tabular}
\end{center}

\question[marks=5, clo=3, bloom=Analyse, topic=Complexity]
The running time of a divide-and-conquer algorithm satisfies
\begin{equation*}
  T(n) = 2\,T\!\left(\frac{n}{2}\right) + \Theta(n), \qquad T(1)=\Theta(1).
\end{equation*}
Solve the recurrence and name an algorithm whose running time it describes.

\question[marks=5, clo=3, bloom=Evaluate, topic=Data structure choice,
          lines=6]
A campus application keeps the 815 enrolled students in memory and looks a
student up by roll number thousands of times a minute, while enrolments are
added only a few times a day. Recommend a data structure and justify the
choice in terms of the operations' costs.

\examsection[20 marks]{Section C: Long questions}

\question[marks=10, clo=3, bloom=Analyse, topic=Sorting]
Consider the array $[38, 27, 43, 3, 9, 82, 10]$.
\qpart[marks=6]
Trace merge sort on it, showing the contents of the array after every merge.
\qpart[marks=4]
Count the comparisons merge sort makes on this input and compare the number
with insertion sort on the same input.

\question[marks=10, clo=4, bloom=Create, topic=Hashing]
A hash table stores the courses a department offers, keyed by course code.
\qpart[marks=4]
Write the insertion routine for separate chaining. One line of code per line
of the answer book:
\begin{center}\small
  \texttt{def insert(table, key, value):}
\end{center}
\qpart[marks=6, space=45mm]
The table has 16 slots and holds 40 courses. Explain what the load factor
does to the cost of a lookup, and describe one way to keep the cost near
constant as the department adds courses. Use the space below for a diagram.

\question[marks=3, bonus, clo=4, bloom=Create, topic=Open question]
State one change to the hash table above that would make its iteration order
predictable, and what it would cost.

\examend

%% ---- For the department's file ---------------------------------------------
%% The distribution table is for the paper setter and the examination office.
%% Delete this page before printing the candidates' copies.
\clearpage
\sausmarkstable

\end{document}

%%==========================================================================
% type = mid | final | sessional | quiz | makeup | supplementary | practical
% solutions = 1 prints the model answers and marks the correct choices
\documentclass[type=final]{saus-exam}

%% ---- The paper -------------------------------------------------------------
\sausdepartment{Department of Computer Science}
\examcourse{CSC-201}{Data Structures and Algorithms}
\examprogramme{BS Computer Science}
\examsemester{III}
\examsession{Spring 2026}
\examdate{16 June 2026}
\examduration{3 hours}
\exammarks{50}
\examcredits{3 (3-0)}
\examinfo{Venue}{Examination Hall 2}

\begin{document}
\makeexamhead

\begin{instructions}
  \item Attempt \textbf{all} questions. Section A is answered on this paper;
        Sections B and C in the answer book provided.
  \item The marks for each question are shown in the margin, with the course
        learning outcome it assesses.
  \item Programmable calculators, mobile telephones and smart watches are not
        allowed in the examination hall.
  \item Write your roll number on every sheet of the answer book.
\end{instructions}

\sausmarksgrid

\examsection[10 marks]{Section A: Multiple choice}

\question[marks=2, clo=1, bloom=Remember, topic=Complexity]
What is the worst-case time complexity of binary search on a sorted array
of $n$ elements?
\begin{choices}[cols=4]
  \choice $O(1)$
  \correctchoice $O(\log n)$
  \choice $O(n)$
  \choice $O(n\log n)$
\end{choices}
\begin{solution}
  Every comparison halves the interval that may still contain the key, so at
  most $\lfloor\log_2 n\rfloor+1$ comparisons are made.
\end{solution}

\question[marks=2, clo=1, bloom=Understand, topic=Stacks and queues]
Which structure serves its items in last-in, first-out order?
\begin{choices}
  \choice Queue
  \correctchoice Stack
  \choice Priority queue
  \choice Circular buffer
\end{choices}

\question[marks=2, clo=1, bloom=Understand, topic=Trees]
A binary search tree holding $n$ keys has height $h$. Searching it takes
time proportional to:
\begin{choices}
  \choice $n$ in every case
  \correctchoice $h$, which is $\log n$ only when the tree is balanced
  \choice $\log n$ in every case
  \choice $n\log n$ in the worst case
\end{choices}

\question[marks=2, clo=2, bloom=Apply, topic=Hashing]
Separate chaining resolves a collision by:
\begin{choices}
  \choice probing the next free slot
  \correctchoice keeping a list of the entries that share a slot
  \choice enlarging the table on every insertion
  \choice rehashing with a second function
\end{choices}

\question[marks=2, clo=2, bloom=Apply, topic=Graphs]
Breadth-first search on an unweighted graph finds:
\begin{choices}
  \choice the minimum spanning tree
  \correctchoice the shortest path in edges from the source
  \choice the strongly connected components
  \choice a topological order
\end{choices}

\examsection[20 marks]{Section B: Short questions}

\question[marks=5, clo=2, bloom=Apply, topic=Recursion]
Write a recursive function that returns the height of a binary tree, and
state its running time in terms of the number of nodes.
\begin{solution}
  The height of an empty tree is $-1$; otherwise it is one more than the
  greater of the heights of the two subtrees. Every node is visited once, so
  the function runs in $O(n)$ time and uses $O(h)$ stack space.
\end{solution}

\question[marks=5, clo=2, bloom=Analyse, topic=Sorting]
Complete the table below, and say which of the three you would choose for
an array that is already nearly sorted.
\begin{center}
  \small
  \begin{tabular}{|l|l|l|l|}
    \hline
    \rowcolor{sausmist}
    \textbf{Algorithm}&\textbf{Best}&\textbf{Average}&\textbf{Worst}\\
    \hline
    Insertion sort&$O(n)$&&\\
    \hline
    Merge sort&&$O(n\log n)$&\\
    \hline
    Quicksort&&&\\
    \hline
  \end{tabular}
\end{center}

\question[marks=5, clo=3, bloom=Analyse, topic=Complexity]
The running time of a divide-and-conquer algorithm satisfies
\begin{equation*}
  T(n) = 2\,T\!\left(\frac{n}{2}\right) + \Theta(n), \qquad T(1)=\Theta(1).
\end{equation*}
Solve the recurrence and name an algorithm whose running time it describes.

\question[marks=5, clo=3, bloom=Evaluate, topic=Data structure choice,
          lines=6]
A campus application keeps the 815 enrolled students in memory and looks a
student up by roll number thousands of times a minute, while enrolments are
added only a few times a day. Recommend a data structure and justify the
choice in terms of the operations' costs.

\examsection[20 marks]{Section C: Long questions}

\question[marks=10, clo=3, bloom=Analyse, topic=Sorting]
Consider the array $[38, 27, 43, 3, 9, 82, 10]$.
\qpart[marks=6]
Trace merge sort on it, showing the contents of the array after every merge.
\qpart[marks=4]
Count the comparisons merge sort makes on this input and compare the number
with insertion sort on the same input.

\question[marks=10, clo=4, bloom=Create, topic=Hashing]
A hash table stores the courses a department offers, keyed by course code.
\qpart[marks=4]
Write the insertion routine for separate chaining. One line of code per line
of the answer book:
\begin{center}\small
  \texttt{def insert(table, key, value):}
\end{center}
\qpart[marks=6, space=45mm]
The table has 16 slots and holds 40 courses. Explain what the load factor
does to the cost of a lookup, and describe one way to keep the cost near
constant as the department adds courses. Use the space below for a diagram.

\question[marks=3, bonus, clo=4, bloom=Create, topic=Open question]
State one change to the hash table above that would make its iteration order
predictable, and what it would cost.

\examend

%% ---- For the department's file ---------------------------------------------
%% The distribution table is for the paper setter and the examination office.
%% Delete this page before printing the candidates' copies.
\clearpage
\sausmarkstable

\end{document}
